% --- Archon physics formalization source begin ---
% archon:physics
% archon:covers PhyXMiniProblems/problem_phyx_mini_0104.lean
% archon:source-report reports/phyx_mini/problem_phyx_mini_0104.source.json

\chapter{Physics problem phyx\_mini\_0104}
\label{ch:PhyXMiniProblems_problem_phyx_mini_0104}

\paragraph{Problem source.}
\#\# Physical scenario

A circularly polarized electromagnetic wave propagating in air has an electric field given by\textbackslash{}[\textbackslash{}vec\{E\} = E \textbackslash{}left[ \textbackslash{}cos(kz - \textbackslash{}omega t) \textbackslash{}hat\{i\} + \textbackslash{}sin(kz - \textbackslash{}omega t) \textbackslash{}hat\{j\} \textbackslash{}right].\textbackslash{}]This wave is incident with an intensity of 150 W/m\$\textasciicircum{}2\$ at the polarizing angle \$\textbackslash{}theta\_p\$ onto a flat interface perpendicular to the \$xz\$-plane with a material that has index of refraction \$n = 1.62\$.The reflected wave is linearly polarized and the refracted wave is elliptically polarized, such that its electric field is characterized as shown in figure. \$e = \textbackslash{}sqrt\{1 - (E\_1 / E\_2)\textasciicircum{}2\}\$.

\#\# Question

Determine the elliptical eccentricity.

\#\# Answer choices

- A: A: 0.449

- B: B: 0.439

- C: C: 0.412

- D: D: 0.521

\#\# Figure caption (auxiliary; use the image as primary evidence)

The image depicts a coordinate system with an ellipse centered at the origin. The axes are labeled \textbackslash{}(x\textbackslash{}) and \textbackslash{}(y\textbackslash{}). The ellipse is drawn in pink, aligned with these axes.

1. **Ellipse**: The ellipse is centered at the origin and stretches diagonally across the \textbackslash{}(x\textbackslash{}) and \textbackslash{}(y\textbackslash{}) axes. 

2. **Axes**: 
   - The horizontal axis (\textbackslash{}(x\textbackslash{})-axis) is labeled with \textbackslash{}(E\_1\textbackslash{}).
   - The vertical axis (\textbackslash{}(y\textbackslash{})-axis) is labeled with \textbackslash{}(E\_2\textbackslash{}).

3. **Vector**: A pink arrow, labeled \textbackslash{}(\textbackslash{}vec\{E\}\textbackslash{}), starts from the origin and points diagonally within the ellipse, indicating direction.

4. **Angle**: There is an angle marked around the origin with a vector labeled \textbackslash{}(\textbackslash{}omega t\textbackslash{}) in black, suggesting rotational or phase angle.

The image likely represents a concept in electromagnetism, depicting an electric field vector rotating in an elliptical path.

\#\# Dataset metadata

Wave Optics; Physical Model Grounding Reasoning, Spatial Relation Reasoning

\paragraph{Recorded answer/context.}
A: A: 0.449

\paragraph{Figure/image path.}
/root/proposal\_for\_physic/science-mango/phyx\_mini\_run/phyx\_data/test\_image/104.png

\paragraph{Formalization target.}
create a compiling Lean file with sorry bodies at `PhyXMiniProblems/problem\_phyx\_mini\_0104.lean`. The Lean declarations must preserve the physical quantities, dimensions or dimensional roles, figure labels, governing-law hypotheses, and final relation expressed by this problem.
Use Mathlib/Physlib names found through LeanExplore where available. If a physics API is missing, introduce faithful local abstractions rather than scalar placeholder aliases.

\begin{theorem}[Physics formalization target]
\label{thm:physics:phyx_mini_0104:target}
\lean{PhyXMiniProblems.ProblemPhyXMini0104.ellipticalEccentricityIsAnswerA}
\uses{def:physics:phyx-mini-0104:phyxminiproblems-problemphyxmini0104-brewsterinterfacesetup, def:physics:phyx-mini-0104:phyxminiproblems-problemphyxmini0104-matchesproblemandfigure, def:physics:phyx-mini-0104:phyxminiproblems-problemphyxmini0104-hasphysicalparameters, def:physics:phyx-mini-0104:phyxminiproblems-problemphyxmini0104-haslinearlypolarizedreflection, def:physics:phyx-mini-0104:phyxminiproblems-problemphyxmini0104-hasellipticallypolarizedtransmission, def:physics:phyx-mini-0104:phyxminiproblems-problemphyxmini0104-isatpolarizingangle, def:physics:phyx-mini-0104:phyxminiproblems-problemphyxmini0104-obeyssnellslaw, def:physics:phyx-mini-0104:phyxminiproblems-problemphyxmini0104-obeysfresneltransmission, def:physics:phyx-mini-0104:phyxminiproblems-problemphyxmini0104-ellipseeccentricity, def:physics:phyx-mini-0104:phyxminiproblems-problemphyxmini0104-isnearestanswerchoice}
The assigned autoformalize agent should translate this physics problem into Lean declarations in the covered file, with theorem and lemma proof bodies written as `by sorry`.
\end{theorem}
\begin{proof}
This is an autoformalization task, not a proof task. Produce statements that can later be proved by the physics prover without weakening the physical model.
\end{proof}

% --- Archon declaration topology begin ---
\section{Declaration topology}
\begin{definition}[Optical Intensity]
  \label{def:physics:phyx-mini-0104:phyxminiproblems-problemphyxmini0104-opticalintensity}
  \lean{PhyXMiniProblems.ProblemPhyXMini0104.OpticalIntensity}
  Optical intensity, with SI readout in watts per square metre
\end{definition}
\begin{proof}
  This is the declared model definition of Optical Intensity.
\end{proof}
\begin{definition}[Electric Field Amplitude]
  \label{def:physics:phyx-mini-0104:phyxminiproblems-problemphyxmini0104-electricfieldamplitude}
  \lean{PhyXMiniProblems.ProblemPhyXMini0104.ElectricFieldAmplitude}
  Peak electric-field strength, with SI readout in volts per metre
\end{definition}
\begin{proof}
  This is the declared model definition of Electric Field Amplitude.
\end{proof}
\begin{definition}[Optical Wave Number]
  \label{def:physics:phyx-mini-0104:phyxminiproblems-problemphyxmini0104-opticalwavenumber}
  \lean{PhyXMiniProblems.ProblemPhyXMini0104.OpticalWaveNumber}
  Wave number, with SI readout in inverse metres
\end{definition}
\begin{proof}
  This is the declared model definition of Optical Wave Number.
\end{proof}
\begin{definition}[Optical Angular Frequency]
  \label{def:physics:phyx-mini-0104:phyxminiproblems-problemphyxmini0104-opticalangularfrequency}
  \lean{PhyXMiniProblems.ProblemPhyXMini0104.OpticalAngularFrequency}
  Angular frequency, with SI readout in radians per second
\end{definition}
\begin{proof}
  This is the declared model definition of Optical Angular Frequency.
\end{proof}
\begin{definition}[intensity In Watts Per Square Meter]
  \label{def:physics:phyx-mini-0104:phyxminiproblems-problemphyxmini0104-intensityinwattspersquaremeter}
  \lean{PhyXMiniProblems.ProblemPhyXMini0104.intensityInWattsPerSquareMeter}
  \uses{def:physics:phyx-mini-0104:phyxminiproblems-problemphyxmini0104-opticalintensity}
  Read an intensity as a number of watts per square metre
\end{definition}
\begin{proof}
  This is the declared model definition of intensity In Watts Per Square Meter.
\end{proof}
\begin{definition}[amplitude In Volts Per Meter]
  \label{def:physics:phyx-mini-0104:phyxminiproblems-problemphyxmini0104-amplitudeinvoltspermeter}
  \lean{PhyXMiniProblems.ProblemPhyXMini0104.amplitudeInVoltsPerMeter}
  \uses{def:physics:phyx-mini-0104:phyxminiproblems-problemphyxmini0104-electricfieldamplitude}
  Read an electric-field amplitude as a number of volts per metre
\end{definition}
\begin{proof}
  This is the declared model definition of amplitude In Volts Per Meter.
\end{proof}
\begin{definition}[wave Number In Inverse Meters]
  \label{def:physics:phyx-mini-0104:phyxminiproblems-problemphyxmini0104-wavenumberininversemeters}
  \lean{PhyXMiniProblems.ProblemPhyXMini0104.waveNumberInInverseMeters}
  \uses{def:physics:phyx-mini-0104:phyxminiproblems-problemphyxmini0104-opticalwavenumber}
  Read a wave number as a number of inverse metres
\end{definition}
\begin{proof}
  This is the declared model definition of wave Number In Inverse Meters.
\end{proof}
\begin{definition}[angular Frequency In Radians Per Second]
  \label{def:physics:phyx-mini-0104:phyxminiproblems-problemphyxmini0104-angularfrequencyinradianspersecond}
  \lean{PhyXMiniProblems.ProblemPhyXMini0104.angularFrequencyInRadiansPerSecond}
  \uses{def:physics:phyx-mini-0104:phyxminiproblems-problemphyxmini0104-opticalangularfrequency}
  Read an angular frequency as a number of radians per second
\end{definition}
\begin{proof}
  This is the declared model definition of angular Frequency In Radians Per Second.
\end{proof}
\begin{definition}[Optical Medium]
  \label{def:physics:phyx-mini-0104:phyxminiproblems-problemphyxmini0104-opticalmedium}
  \lean{PhyXMiniProblems.ProblemPhyXMini0104.OpticalMedium}
  The homogeneous media on the two sides of the interface
\end{definition}
\begin{proof}
  This is the declared model definition of Optical Medium.
\end{proof}
\begin{definition}[Wave Kind]
  \label{def:physics:phyx-mini-0104:phyxminiproblems-problemphyxmini0104-wavekind}
  \lean{PhyXMiniProblems.ProblemPhyXMini0104.WaveKind}
  The three waves associated with the interface
\end{definition}
\begin{proof}
  This is the declared model definition of Wave Kind.
\end{proof}
\begin{definition}[medium]
  \label{def:physics:phyx-mini-0104:phyxminiproblems-problemphyxmini0104-wavekind-medium}
  \lean{PhyXMiniProblems.ProblemPhyXMini0104.WaveKind.medium}
  \uses{def:physics:phyx-mini-0104:phyxminiproblems-problemphyxmini0104-opticalmedium, def:physics:phyx-mini-0104:phyxminiproblems-problemphyxmini0104-wavekind}
  Medium occupied by each wave in the idealized interface model
\end{definition}
\begin{proof}
  This is the declared model definition of medium.
\end{proof}
\begin{definition}[Polarization Component]
  \label{def:physics:phyx-mini-0104:phyxminiproblems-problemphyxmini0104-polarizationcomponent}
  \lean{PhyXMiniProblems.ProblemPhyXMini0104.PolarizationComponent}
  The usual p and s components relative to the plane of incidence
\end{definition}
\begin{proof}
  This is the declared model definition of Polarization Component.
\end{proof}
\begin{definition}[Cartesian Axis]
  \label{def:physics:phyx-mini-0104:phyxminiproblems-problemphyxmini0104-cartesianaxis}
  \lean{PhyXMiniProblems.ProblemPhyXMini0104.CartesianAxis}
  Cartesian axes named in the incident-wave formula and ellipse figure
\end{definition}
\begin{proof}
  This is the declared model definition of Cartesian Axis.
\end{proof}
\begin{definition}[Coordinate Plane]
  \label{def:physics:phyx-mini-0104:phyxminiproblems-problemphyxmini0104-coordinateplane}
  \lean{PhyXMiniProblems.ProblemPhyXMini0104.CoordinatePlane}
  Coordinate planes used to record the source's interface orientation
\end{definition}
\begin{proof}
  This is the declared model definition of Coordinate Plane.
\end{proof}
\begin{definition}[Figure Amplitude Label]
  \label{def:physics:phyx-mini-0104:phyxminiproblems-problemphyxmini0104-figureamplitudelabel}
  \lean{PhyXMiniProblems.ProblemPhyXMini0104.FigureAmplitudeLabel}
  The two semiaxis labels printed in the primary ellipse image
\end{definition}
\begin{proof}
  This is the declared model definition of Figure Amplitude Label.
\end{proof}
\begin{definition}[axis]
  \label{def:physics:phyx-mini-0104:phyxminiproblems-problemphyxmini0104-figureamplitudelabel-axis}
  \lean{PhyXMiniProblems.ProblemPhyXMini0104.FigureAmplitudeLabel.axis}
  \uses{def:physics:phyx-mini-0104:phyxminiproblems-problemphyxmini0104-cartesianaxis, def:physics:phyx-mini-0104:phyxminiproblems-problemphyxmini0104-figureamplitudelabel}
  The image places E₁ on the horizontal axis and E₂ on the vertical axis
\end{definition}
\begin{proof}
  This is the declared model definition of axis.
\end{proof}
\begin{definition}[component]
  \label{def:physics:phyx-mini-0104:phyxminiproblems-problemphyxmini0104-figureamplitudelabel-component}
  \lean{PhyXMiniProblems.ProblemPhyXMini0104.FigureAmplitudeLabel.component}
  \uses{def:physics:phyx-mini-0104:phyxminiproblems-problemphyxmini0104-polarizationcomponent, def:physics:phyx-mini-0104:phyxminiproblems-problemphyxmini0104-figureamplitudelabel}
  Polarization component represented by each principal ellipse axis
\end{definition}
\begin{proof}
  This is the declared model definition of component.
\end{proof}
\begin{definition}[Brewster Interface Setup]
  \label{def:physics:phyx-mini-0104:phyxminiproblems-problemphyxmini0104-brewsterinterfacesetup}
  \lean{PhyXMiniProblems.ProblemPhyXMini0104.BrewsterInterfaceSetup}
  \uses{def:physics:phyx-mini-0104:phyxminiproblems-problemphyxmini0104-opticalintensity, def:physics:phyx-mini-0104:phyxminiproblems-problemphyxmini0104-electricfieldamplitude, def:physics:phyx-mini-0104:phyxminiproblems-problemphyxmini0104-opticalwavenumber, def:physics:phyx-mini-0104:phyxminiproblems-problemphyxmini0104-opticalangularfrequency, def:physics:phyx-mini-0104:phyxminiproblems-problemphyxmini0104-opticalmedium, def:physics:phyx-mini-0104:phyxminiproblems-problemphyxmini0104-wavekind, def:physics:phyx-mini-0104:phyxminiproblems-problemphyxmini0104-polarizationcomponent, def:physics:phyx-mini-0104:phyxminiproblems-problemphyxmini0104-coordinateplane}
  Physical quantities in the Brewster-angle interface experiment. For each wave, electricField is expressed in a local orthonormal frame whose first two coordinates are the parallel and perpendicular transverse polarization directions and whose third coordinate is the propagation direction. The scalar component values are SI electric-field readouts
\end{definition}
\begin{proof}
  This is the declared model definition of Brewster Interface Setup.
\end{proof}
\begin{definition}[figure Semi Axis Amplitude]
  \label{def:physics:phyx-mini-0104:phyxminiproblems-problemphyxmini0104-figuresemiaxisamplitude}
  \lean{PhyXMiniProblems.ProblemPhyXMini0104.figureSemiAxisAmplitude}
  \uses{def:physics:phyx-mini-0104:phyxminiproblems-problemphyxmini0104-electricfieldamplitude, def:physics:phyx-mini-0104:phyxminiproblems-problemphyxmini0104-figureamplitudelabel, def:physics:phyx-mini-0104:phyxminiproblems-problemphyxmini0104-brewsterinterfacesetup}
  The physical amplitude represented by a label in the ellipse image
\end{definition}
\begin{proof}
  This is the declared model definition of figure Semi Axis Amplitude.
\end{proof}
\begin{definition}[Has Quadrature Polarization]
  \label{def:physics:phyx-mini-0104:phyxminiproblems-problemphyxmini0104-hasquadraturepolarization}
  \lean{PhyXMiniProblems.ProblemPhyXMini0104.HasQuadraturePolarization}
  \uses{def:physics:phyx-mini-0104:phyxminiproblems-problemphyxmini0104-amplitudeinvoltspermeter, def:physics:phyx-mini-0104:phyxminiproblems-problemphyxmini0104-wavekind, def:physics:phyx-mini-0104:phyxminiproblems-problemphyxmini0104-brewsterinterfacesetup}
  At every spacetime point, the two transverse components are in quadrature and the longitudinal electric-field component vanishes
\end{definition}
\begin{proof}
  This is the declared model definition of Has Quadrature Polarization.
\end{proof}
\begin{definition}[Matches Problem And Figure]
  \label{def:physics:phyx-mini-0104:phyxminiproblems-problemphyxmini0104-matchesproblemandfigure}
  \lean{PhyXMiniProblems.ProblemPhyXMini0104.MatchesProblemAndFigure}
  \uses{def:physics:phyx-mini-0104:phyxminiproblems-problemphyxmini0104-intensityinwattspersquaremeter, def:physics:phyx-mini-0104:phyxminiproblems-problemphyxmini0104-amplitudeinvoltspermeter, def:physics:phyx-mini-0104:phyxminiproblems-problemphyxmini0104-wavenumberininversemeters, def:physics:phyx-mini-0104:phyxminiproblems-problemphyxmini0104-angularfrequencyinradianspersecond, def:physics:phyx-mini-0104:phyxminiproblems-problemphyxmini0104-brewsterinterfacesetup, def:physics:phyx-mini-0104:phyxminiproblems-problemphyxmini0104-figuresemiaxisamplitude, def:physics:phyx-mini-0104:phyxminiproblems-problemphyxmini0104-hasquadraturepolarization}
  Problem-text and primary-image readouts. This records the 150 W/m² intensity, index 1.62, incident phase kz - ωt, equal incident components, the stated xz orientation, and the image's ordering of the two semiaxes. It does not record an eccentricity or an answer choice
\end{definition}
\begin{proof}
  This is the declared model definition of Matches Problem And Figure.
\end{proof}
\begin{definition}[Is Acute Optical Angle]
  \label{def:physics:phyx-mini-0104:phyxminiproblems-problemphyxmini0104-isacuteopticalangle}
  \lean{PhyXMiniProblems.ProblemPhyXMini0104.IsAcuteOpticalAngle}
  An angle is on the acute branch used by the physical interface diagram
\end{definition}
\begin{proof}
  This is the declared model definition of Is Acute Optical Angle.
\end{proof}
\begin{definition}[Has Physical Parameters]
  \label{def:physics:phyx-mini-0104:phyxminiproblems-problemphyxmini0104-hasphysicalparameters}
  \lean{PhyXMiniProblems.ProblemPhyXMini0104.HasPhysicalParameters}
  \uses{def:physics:phyx-mini-0104:phyxminiproblems-problemphyxmini0104-intensityinwattspersquaremeter, def:physics:phyx-mini-0104:phyxminiproblems-problemphyxmini0104-amplitudeinvoltspermeter, def:physics:phyx-mini-0104:phyxminiproblems-problemphyxmini0104-wavenumberininversemeters, def:physics:phyx-mini-0104:phyxminiproblems-problemphyxmini0104-angularfrequencyinradianspersecond, def:physics:phyx-mini-0104:phyxminiproblems-problemphyxmini0104-brewsterinterfacesetup, def:physics:phyx-mini-0104:phyxminiproblems-problemphyxmini0104-isacuteopticalangle}
  Positivity and branch conditions for the optical model
\end{definition}
\begin{proof}
  This is the declared model definition of Has Physical Parameters.
\end{proof}
\begin{definition}[Has Linearly Polarized Reflection]
  \label{def:physics:phyx-mini-0104:phyxminiproblems-problemphyxmini0104-haslinearlypolarizedreflection}
  \lean{PhyXMiniProblems.ProblemPhyXMini0104.HasLinearlyPolarizedReflection}
  \uses{def:physics:phyx-mini-0104:phyxminiproblems-problemphyxmini0104-amplitudeinvoltspermeter, def:physics:phyx-mini-0104:phyxminiproblems-problemphyxmini0104-brewsterinterfacesetup}
  The reflected wave is s-polarized: its p component vanishes
\end{definition}
\begin{proof}
  This is the declared model definition of Has Linearly Polarized Reflection.
\end{proof}
\begin{definition}[Has Elliptically Polarized Transmission]
  \label{def:physics:phyx-mini-0104:phyxminiproblems-problemphyxmini0104-hasellipticallypolarizedtransmission}
  \lean{PhyXMiniProblems.ProblemPhyXMini0104.HasEllipticallyPolarizedTransmission}
  \uses{def:physics:phyx-mini-0104:phyxminiproblems-problemphyxmini0104-brewsterinterfacesetup, def:physics:phyx-mini-0104:phyxminiproblems-problemphyxmini0104-figuresemiaxisamplitude, def:physics:phyx-mini-0104:phyxminiproblems-problemphyxmini0104-hasquadraturepolarization}
  The refracted wave follows the ellipse and has unequal principal axes
\end{definition}
\begin{proof}
  This is the declared model definition of Has Elliptically Polarized Transmission.
\end{proof}
\begin{definition}[angle From Degrees]
  \label{def:physics:phyx-mini-0104:phyxminiproblems-problemphyxmini0104-anglefromdegrees}
  \lean{PhyXMiniProblems.ProblemPhyXMini0104.angleFromDegrees}
  Convert a numerical degree readout to Mathlib's physical angle type
\end{definition}
\begin{proof}
  This is the declared model definition of angle From Degrees.
\end{proof}
\begin{definition}[Is At Polarizing Angle]
  \label{def:physics:phyx-mini-0104:phyxminiproblems-problemphyxmini0104-isatpolarizingangle}
  \lean{PhyXMiniProblems.ProblemPhyXMini0104.IsAtPolarizingAngle}
  \uses{def:physics:phyx-mini-0104:phyxminiproblems-problemphyxmini0104-brewsterinterfacesetup, def:physics:phyx-mini-0104:phyxminiproblems-problemphyxmini0104-anglefromdegrees}
  Characterization of incidence at the polarizing angle: the two ray angles are complementary and the tangent of the incident angle is the index ratio
\end{definition}
\begin{proof}
  This is the declared model definition of Is At Polarizing Angle.
\end{proof}
\begin{definition}[Obeys Snells Law]
  \label{def:physics:phyx-mini-0104:phyxminiproblems-problemphyxmini0104-obeyssnellslaw}
  \lean{PhyXMiniProblems.ProblemPhyXMini0104.ObeysSnellsLaw}
  \uses{def:physics:phyx-mini-0104:phyxminiproblems-problemphyxmini0104-brewsterinterfacesetup}
  Snell's law at the air-material interface
\end{definition}
\begin{proof}
  This is the declared model definition of Obeys Snells Law.
\end{proof}
\begin{definition}[fresnel Transmission Coefficient]
  \label{def:physics:phyx-mini-0104:phyxminiproblems-problemphyxmini0104-fresneltransmissioncoefficient}
  \lean{PhyXMiniProblems.ProblemPhyXMini0104.fresnelTransmissionCoefficient}
  \uses{def:physics:phyx-mini-0104:phyxminiproblems-problemphyxmini0104-polarizationcomponent, def:physics:phyx-mini-0104:phyxminiproblems-problemphyxmini0104-brewsterinterfacesetup}
  Electric-field Fresnel transmission coefficient for a nonmagnetic dielectric interface, separately for p and s polarization
\end{definition}
\begin{proof}
  This is the declared model definition of fresnel Transmission Coefficient.
\end{proof}
\begin{definition}[Obeys Fresnel Transmission]
  \label{def:physics:phyx-mini-0104:phyxminiproblems-problemphyxmini0104-obeysfresneltransmission}
  \lean{PhyXMiniProblems.ProblemPhyXMini0104.ObeysFresnelTransmission}
  \uses{def:physics:phyx-mini-0104:phyxminiproblems-problemphyxmini0104-amplitudeinvoltspermeter, def:physics:phyx-mini-0104:phyxminiproblems-problemphyxmini0104-brewsterinterfacesetup, def:physics:phyx-mini-0104:phyxminiproblems-problemphyxmini0104-fresneltransmissioncoefficient}
  Fresnel transmission scales each incident component independently
\end{definition}
\begin{proof}
  This is the declared model definition of Obeys Fresnel Transmission.
\end{proof}
\begin{definition}[ellipse Coordinates In Volts Per Meter]
  \label{def:physics:phyx-mini-0104:phyxminiproblems-problemphyxmini0104-ellipsecoordinatesinvoltspermeter}
  \lean{PhyXMiniProblems.ProblemPhyXMini0104.ellipseCoordinatesInVoltsPerMeter}
  \uses{def:physics:phyx-mini-0104:phyxminiproblems-problemphyxmini0104-amplitudeinvoltspermeter, def:physics:phyx-mini-0104:phyxminiproblems-problemphyxmini0104-brewsterinterfacesetup, def:physics:phyx-mini-0104:phyxminiproblems-problemphyxmini0104-figuresemiaxisamplitude}
  The rotating field-vector coordinates depicted at phase ωt
\end{definition}
\begin{proof}
  This is the declared model definition of ellipse Coordinates In Volts Per Meter.
\end{proof}
\begin{definition}[ellipse Eccentricity]
  \label{def:physics:phyx-mini-0104:phyxminiproblems-problemphyxmini0104-ellipseeccentricity}
  \lean{PhyXMiniProblems.ProblemPhyXMini0104.ellipseEccentricity}
  \uses{def:physics:phyx-mini-0104:phyxminiproblems-problemphyxmini0104-amplitudeinvoltspermeter, def:physics:phyx-mini-0104:phyxminiproblems-problemphyxmini0104-brewsterinterfacesetup, def:physics:phyx-mini-0104:phyxminiproblems-problemphyxmini0104-figuresemiaxisamplitude}
  Standard eccentricity of the figure's E₁-major, E₂-minor ellipse
\end{definition}
\begin{proof}
  This is the declared model definition of ellipse Eccentricity.
\end{proof}
\begin{definition}[Answer Choice]
  \label{def:physics:phyx-mini-0104:phyxminiproblems-problemphyxmini0104-answerchoice}
  \lean{PhyXMiniProblems.ProblemPhyXMini0104.AnswerChoice}
  Labels of the four displayed answer choices
\end{definition}
\begin{proof}
  This is the declared model definition of Answer Choice.
\end{proof}
\begin{definition}[eccentricity Readout]
  \label{def:physics:phyx-mini-0104:phyxminiproblems-problemphyxmini0104-answerchoice-eccentricityreadout}
  \lean{PhyXMiniProblems.ProblemPhyXMini0104.AnswerChoice.eccentricityReadout}
  \uses{def:physics:phyx-mini-0104:phyxminiproblems-problemphyxmini0104-answerchoice}
  Dimensionless eccentricity printed beside each answer label
\end{definition}
\begin{proof}
  This is the declared model definition of eccentricity Readout.
\end{proof}
\begin{definition}[recorded Answer Choice]
  \label{def:physics:phyx-mini-0104:phyxminiproblems-problemphyxmini0104-recordedanswerchoice}
  \lean{PhyXMiniProblems.ProblemPhyXMini0104.recordedAnswerChoice}
  \uses{def:physics:phyx-mini-0104:phyxminiproblems-problemphyxmini0104-answerchoice}
  Dataset metadata: the recorded answer is A; this is not a theorem premise
\end{definition}
\begin{proof}
  This is the declared model definition of recorded Answer Choice.
\end{proof}
\begin{definition}[Is Nearest Answer Choice]
  \label{def:physics:phyx-mini-0104:phyxminiproblems-problemphyxmini0104-isnearestanswerchoice}
  \lean{PhyXMiniProblems.ProblemPhyXMini0104.IsNearestAnswerChoice}
  \uses{def:physics:phyx-mini-0104:phyxminiproblems-problemphyxmini0104-answerchoice}
  A displayed choice is strictly nearest to the physical eccentricity
\end{definition}
\begin{proof}
  This is the declared model definition of Is Nearest Answer Choice.
\end{proof}
\begin{lemma}[transmitted Axis Ratio At Polarizing Angle]
  \label{lem:physics:phyx-mini-0104:phyxminiproblems-problemphyxmini0104-transmittedaxisratioatpolarizingangle}
  \lean{PhyXMiniProblems.ProblemPhyXMini0104.transmittedAxisRatioAtPolarizingAngle}
  \uses{def:physics:phyx-mini-0104:phyxminiproblems-problemphyxmini0104-amplitudeinvoltspermeter, def:physics:phyx-mini-0104:phyxminiproblems-problemphyxmini0104-brewsterinterfacesetup, def:physics:phyx-mini-0104:phyxminiproblems-problemphyxmini0104-figuresemiaxisamplitude, def:physics:phyx-mini-0104:phyxminiproblems-problemphyxmini0104-matchesproblemandfigure, def:physics:phyx-mini-0104:phyxminiproblems-problemphyxmini0104-hasphysicalparameters, def:physics:phyx-mini-0104:phyxminiproblems-problemphyxmini0104-isatpolarizingangle, def:physics:phyx-mini-0104:phyxminiproblems-problemphyxmini0104-obeyssnellslaw, def:physics:phyx-mini-0104:phyxminiproblems-problemphyxmini0104-obeysfresneltransmission}
  At Brewster incidence the Fresnel laws make the minor-to-major semiaxis ratio 2 n₁ n₂ / (n₁² + n₂²)
\end{lemma}
\begin{proof}
  The conclusion follows from the governing relations and figure readouts named in the statement of transmitted Axis Ratio At Polarizing Angle.
\end{proof}
\begin{lemma}[eccentricity From Fresnel Transmission]
  \label{lem:physics:phyx-mini-0104:phyxminiproblems-problemphyxmini0104-eccentricityfromfresneltransmission}
  \lean{PhyXMiniProblems.ProblemPhyXMini0104.eccentricityFromFresnelTransmission}
  \uses{def:physics:phyx-mini-0104:phyxminiproblems-problemphyxmini0104-brewsterinterfacesetup, def:physics:phyx-mini-0104:phyxminiproblems-problemphyxmini0104-matchesproblemandfigure, def:physics:phyx-mini-0104:phyxminiproblems-problemphyxmini0104-hasphysicalparameters, def:physics:phyx-mini-0104:phyxminiproblems-problemphyxmini0104-isatpolarizingangle, def:physics:phyx-mini-0104:phyxminiproblems-problemphyxmini0104-obeyssnellslaw, def:physics:phyx-mini-0104:phyxminiproblems-problemphyxmini0104-obeysfresneltransmission, def:physics:phyx-mini-0104:phyxminiproblems-problemphyxmini0104-ellipseeccentricity}
  The corresponding exact Fresnel expression for the ellipse eccentricity
\end{lemma}
\begin{proof}
  The conclusion follows from the governing relations and figure readouts named in the statement of eccentricity From Fresnel Transmission.
\end{proof}
% --- Archon declaration topology end ---
% --- Archon physics formalization source end ---
